\chapter{Quotient Negative-Inversion Zero-Set No-Go}
\label{chap:quotient-negative-inversion-zero-set-no-go}

\status{Exact quotient-level no-go theorem proved. Only finitely many roots of $F$ can be paired with another root by $w\mapsto1/(16w)$. SOH-G003 and RH remain open.}

Chapter~\ref{chap:negative-inversion-zero-set-no-go} proves that the canonical
negative inversion cannot globally preserve the xi zero set.  The present
chapter transfers that firewall to the even quotient entire function
\[
 \xi\!\left(\frac12+z\right)=F(z^2),
\]
where the G012 operator takes its simplest reciprocal form.

\section{The even quotient and its positive center}

By the evenness of the centered xi function there is a unique entire function
$F$ satisfying
\begin{equation}
 \xi\!\left(\frac12+z\right)=F(z^2).
 \label{eq:g015-F-def}
\end{equation}
Chapter~\ref{chap:positive-riemann-kernel} proves the coefficient expansion
\begin{equation}
 F(w)=\sum_{k\ge0}a_k w^k,
 \qquad a_k>0.
 \label{eq:g015-positive-coefficients}
\end{equation}
Consequently
\begin{equation}
 \boxed{F(0)=\xi(1/2)>0}
 \label{eq:g015-F0-positive}
\end{equation}
and, more generally,
\begin{equation}
 \boxed{F(x)>0\quad(x\ge0).}
 \label{eq:g015-positive-axis}
\end{equation}

\section{The quotient negative inversion}

In the centered coordinate $z=s-1/2$, G012 gives
\[
 N_z(z)=-\frac1{4z}.
\]
With
\[
 w=z^2
\]
this descends to
\begin{equation}
 \boxed{
 J(w)=\frac1{16w}.
 }
 \label{eq:g015-J}
\end{equation}
The quotient map is an involution,
\[
 J(J(w))=w,
\]
and its fixed points are
\begin{equation}
 \boxed{w=\pm\frac14.}
 \label{eq:g015-fixed-values}
\end{equation}
Equation~\eqref{eq:g015-positive-axis} immediately excludes one of them from
the root set:
\begin{equation}
 \boxed{F(1/4)>0.}
 \label{eq:g015-plus-quarter-nonroot}
\end{equation}
No claim about $F(-1/4)$ is made at this stage.

\section{The paired quotient-root set}

Let
\[
 Z_F=\{w\in\C:F(w)=0\}
\]
and define
\begin{equation}
 P_J
 =\{w\in Z_F:F(J(w))=0\}.
 \label{eq:g015-paired-set}
\end{equation}
The set $Z_F$ is infinite.  Indeed, Hardy's theorem provides infinitely many
critical-line zeros \citep{hardy1914}; if
\[
 \rho_n=\frac12+i\gamma_n,
 \qquad |\gamma_n|\to\infty,
\]
then \cref{eq:g015-F-def} gives the distinct quotient roots
\[
 w_n=(\rho_n-1/2)^2=-\gamma_n^2.
\]

\begin{theorem}[Quotient paired-root no-go]
\label{thm:g015-paired-finite}
The paired quotient-root set $P_J$ is finite.
\end{theorem}

\begin{proof}
Assume that $P_J$ is infinite.  Since $F$ is a non-zero entire function, its
zeros are isolated, so an infinite subset of its zero set must be unbounded.
Choose
\[
 w_n\in P_J,
 \qquad |w_n|\to\infty.
\]
By definition,
\[
 F(J(w_n))=0
\]
for every $n$.  But \cref{eq:g015-J} gives
\[
 J(w_n)=\frac1{16w_n}\longrightarrow0.
\]
Continuity of $F$ therefore forces
\[
 F(0)
 =\lim_{n\to\infty}F(J(w_n))
 =0,
\]
contradicting \cref{eq:g015-F0-positive}.  Hence $P_J$ is finite.
\end{proof}

\begin{corollary}[No global quotient-root invariance]
\label{cor:g015-no-global-invariance}
The complete root set of $F$ is not invariant under $J$:
\[
 \boxed{J(Z_F)\ne Z_F.}
\]
Moreover, all but finitely many roots of $F$ are sent by $J$ to non-roots.
\end{corollary}

\section{Exact orbit classification of the exceptional set}

The paired set is itself $J$-invariant.  Indeed, if $w\in P_J$, then both
$w$ and $J(w)$ are roots and
\[
 J(J(w))=w.
\]
Since $P_J$ is finite and $J$ is an involution, it decomposes into orbits of
size one or two.

The size-one orbits are fixed values of $J$, so by
\cref{eq:g015-fixed-values} they lie in
\[
 \left\{\frac14,-\frac14\right\}.
\]
The positive fixed value is excluded by
\cref{eq:g015-plus-quarter-nonroot}.  Therefore:

\begin{theorem}[Exceptional-orbit classification]
\label{thm:g015-orbit-classification}
The finite set $P_J$ is a disjoint union of two-cycles, possibly together with
the single fixed point $-1/4$ if $F(-1/4)=0$.
\end{theorem}

Equivalently,
\[
 \boxed{
 P_J
 =\Bigl(\bigsqcup_{m=1}^{M}\{w_m,J(w_m)\}\Bigr)
 \sqcup E,
 }
\]
where
\[
 E\in\{\varnothing,\{-1/4\}\}.
\]
Thus the parity of $|P_J|$ is controlled entirely by the unresolved value
$F(-1/4)$: it is even unless $-1/4$ is itself a root.

\section{Consequences for the real-rootedness frontier}

The real-rootedness target of G003 asks whether every zero of $F$ lies on the
non-positive real axis.  G015 neither proves nor disproves that statement.
It removes a different possible shortcut.

Even if RH is true and all roots of $F$ are negative real, the map
\[
 J(w)=\frac1{16w}
\]
would preserve the negative real axis merely as a geometric locus.  By
\cref{thm:g015-paired-finite}, it still cannot globally permute the roots of
$F$.

Hence reciprocal quotient-root invariance is not a missing condition for
G003 and must not be inserted into the proof programme as though it followed
from RH or from the inverse-boundary geometry.

\section{Relation to the previous operator layers}

G012 establishes the exact projective operator and its quotient action.
G013 supplies the Pauli/SU(2) spinor lift.  G014 proves the corresponding
$s$-plane spectral no-go.  G015 closes the same loophole directly in the
entire-function quotient used by the real-rootedness reduction:
\[
 \boxed{
 \text{quotient operator symmetry}
 \not\Rightarrow
 \text{quotient-root symmetry}.
 }
\]
The operator geometry remains exact.  The rejected statement is only global
spectral invariance.

\section{Numerical regression is not the proof}

The accompanying test and receipt map the first few non-trivial xi zeros to
quotient roots
\[
 w_n=-\gamma_n^2
\]
and verify that their images $J(w_n)$ lie close to zero while
$F(J(w_n))$ remains far from zero.  The same machinery checks branch
independence of the numerical $F$ evaluator and the two quotient fixed values.
These calculations are implementation diagnostics only.  The theorem is the
analytic compactness/continuity contradiction above.

\section{Proof firewall}

Proved in SOH-G015:
\begin{itemize}
 \item the exact quotient involution $J(w)=1/(16w)$;
 \item the fixed values $w=\pm1/4$;
 \item $F(0)>0$ and $F(1/4)>0$ from the previously proved coefficient positivity;
 \item finiteness of the paired quotient-root set $P_J$;
 \item failure of global $F$-root invariance under $J$;
 \item cofinite failure of quotient-root pairing;
 \item decomposition of $P_J$ into two-cycles and possibly the fixed root $-1/4$.
\end{itemize}

Not proved or claimed in SOH-G015:
\begin{itemize}
 \item whether $F(-1/4)$ vanishes;
 \item real-rootedness of $F$;
 \item PF$_\infty$;
 \item SOH-G003;
 \item the Riemann Hypothesis.
\end{itemize}
